\chapter{A Concrete Naming Certificate for $\CofinalCompactIntervalUp$}
\label{ch:concrete-instances-cofinal-compact-interval-namecert}
\origin{human}

Following Bishop and Bridges constructive compact-interval arguments, the $\CofinalCompactIntervalUp$ packet names the finite cofinal window by which a closed bounded interval request is refined to displayed locator, dyadic, stream, regular-rational, and Real-facing rows. It sits after the closed bounded interval source of \autoref{ch:concrete-instances-closedboundedinterval-namecert}, the compact interval locator of \autoref{ch:concrete-instances-compact-interval-locator-namecert}, the dyadic tolerance ledger of \autoref{ch:concrete-instances-dyadicratcore-namecert}, the finite observation schedule of \autoref{ch:concrete-instances-streamname-namecert}, the regular rational readback of \autoref{ch:concrete-instances-regseqrat-namecert}, and the terminal real-name seal of \autoref{ch:concrete-instances-real-namecert}. The packet gives downstream interval compactness, Cousin-style finite refinement, and Heine--Borel interval consumers a named finite-window route without importing open-cover compactness, selected subcovers, quotient real equality, countable choice, trichotomy, or ambient $\RealUp$ completeness.

\begin{definition}[Cofinal compact interval carrier]
\label{def:cofinal-compact-interval-carrier}
A $\CofinalCompactIntervalUp$ carrier is a finite $\BHist$ packet
$$
K=(I,L,D,W,R,E,H,C,P,N)
$$
with the following displayed rows:
\begin{enumerate}
\item $I$ is the $\ClosedBoundedIntervalUp$ source row, exposing lower and upper located endpoints, endpoint order, rational containment, dyadic containment, and the base interval naming surface.
\item $L$ is the $\CompactIntervalLocatorUp$ row, recording the located comparison and tolerance-to-window service used to refine interval requests.
\item $D$ is the $\DyadicRatCoreUp$ ledger carrying the dyadic tolerances, endpoint cells, and finite mesh comparisons demanded by $I$ and $L$.
\item $W$ is the $\StreamNameUp$ finite-window row opened at the cofinal dyadic depth chosen for the displayed request.
\item $R$ is the $\RegSeqRatUp$ regular rational readback attached to the same interval, locator, dyadic, and stream rows.
\item $E$ is the $\RealUp$ seal row reached only after $I,L,D,W,R$ have been displayed.
\item $H$ records componentwise $\hsame$ transport for the interval source, locator, dyadic ledger, stream windows, regular readback, real seal, replay, provenance, and local naming rows.
\item $C$ records displayed $\Cont$ replay from a compact-interval request through $I,L,D,W,R,E,H,P,N$.
\item $P$ is the $\Pkg$ provenance over $\CofinalCompactIntervalUp$, $\ClosedBoundedIntervalUp$, $\CompactIntervalLocatorUp$, $\DyadicRatCoreUp$, $\StreamNameUp$, $\RegSeqRatUp$, $\RealUp$, $\BHist$, $\hsame$, $\Cont$, and $\NameCert$.
\item $N$ is the local $\NameCert_{\CofinalCompactIntervalUp}$ row.
\end{enumerate}
The classifier admits only packets with these ten rows. No coordinate stores open-cover compactness, a selected subcover, quotient real equality, countable choice, trichotomy, host interval equality, or an ambient completeness theorem for $\RealUp$.
\end{definition}

\begin{theorem}[Cofinal compact interval NameCert obligations]
\label{thm:cofinal-compact-interval-namecert-obligations}
For every accepted $\CofinalCompactIntervalUp$ carrier $K=(I,L,D,W,R,E,H,C,P,N)$, the local $\NameCert_{\CofinalCompactIntervalUp}$ surface consists exactly of the following obligation rows:
\begin{enumerate}
\item carrier admission for the closed bounded interval source, compact locator, dyadic tolerance ledger, stream windows, regular rational readback, real seal, transports, continuations, provenance, and local naming;
\item interval-source exactness, requiring every endpoint, containment, and interval-local request to enter through $I$ before the locator row is read;
\item locator cofinality, requiring $L$ and $D$ to convert any displayed compact-interval request into a finite dyadic window deep enough for the requested observation;
\item real-handoff ordering, requiring every $\RealUp$ read to pass through $I,L,D,W,R$ before $E,H,C,P,N$ are exposed;
\item ledger non-escape, excluding open-cover compactness, selected subcovers, quotient real equality, countable choice, trichotomy, host interval equality, and ambient $\RealUp$ completeness.
\end{enumerate}
\end{theorem}
\begin{proof}
Project the accepted carrier of \autoref{def:cofinal-compact-interval-carrier}. The local naming surface has exactly the displayed rows $I,L,D,W,R,E,H,C,P,N$.

The interval-source exactness row is supplied by reading $I$ through \autoref{ch:concrete-instances-closedboundedinterval-namecert}. It exposes the lower endpoint, upper endpoint, endpoint-order, rational containment, dyadic containment, finite-window, regular-readback, and Real-seal boundary as displayed interval data.

The locator cofinality row is then read through \autoref{ch:concrete-instances-compact-interval-locator-namecert} and the dyadic tolerance ledger of \autoref{ch:concrete-instances-dyadicratcore-namecert}. These rows force the selected finite window to be a displayed refinement of the same interval request before $W$ or $R$ can be consumed.

Finally the real-handoff ordering row sends every Real-facing read through $W$ and $R$ before $E$ and the structural rows $H,C,P,N$ are exposed. Since no coordinate of the carrier is an open cover, selected subcover, quotient-real equality witness, countable-choice schedule, trichotomy row, host interval equality, or ambient completeness theorem, the ledger non-escape obligation is exactly the absence of those rows.
\end{proof}

\begin{theorem}[Cofinal compact interval Real handoff]
\label{thm:cofinal-compact-interval-real-handoff}
Every $\RealUp$, interval-compactness, Cousin-refinement, or Heine--Borel-style consumer admitted by an accepted $\CofinalCompactIntervalUp$ packet reads the compact interval request through
$$
\ClosedBoundedIntervalUp,\ \CompactIntervalLocatorUp,\ \DyadicRatCoreUp,\ \StreamNameUp,\ \RegSeqRatUp,\ \RealUp
$$
in that order. The handoff exposes a finite cofinal window and its real seal; it exposes no open-cover compactness row, selected subcover, quotient real equality, countable-choice selector, trichotomy decision, host interval equality, or ambient $\RealUp$ completeness witness.
\end{theorem}
\begin{proof}
Start with an admitted consumer. By \autoref{thm:cofinal-compact-interval-namecert-obligations}, the read first enters through the closed bounded interval row $I$, so the endpoint and containment data are fixed before any compact-window refinement is available.

Next the same obligation surface requires locator cofinality. The locator row $L$ and dyadic ledger $D$ therefore choose a displayed finite dyadic window for the same request before the stream row $W$ and regular readback $R$ are visible.

The real-handoff ordering row then sends the read through $R$ to the real seal $E$, followed only by transport, replay, provenance, and the local naming row. Because the carrier has no other coordinate, the displayed six-object route exhausts the consumer-facing handoff and excludes the listed compactness, choice, quotient, order, host, and completeness data.
\end{proof}

\falsifiablePrediction{An accepted $\CofinalCompactIntervalUp$ consumer that reaches $\RealUp$ without first reading the closed bounded interval row, compact interval locator, dyadic tolerance ledger, finite stream window, and regular rational readback refutes the carrier surface.}

\independenceWitness{closedboundedinterval, compact-interval-locator, dyadicratcore, streamname, regseqrat, real: none is bijective to $\CofinalCompactIntervalUp$ because this packet jointly carries the interval source, compact locator, cofinal dyadic window, stream schedule, regular readback, Real seal, transport, replay, provenance, and local naming.}

\closureat{CofinalCompactIntervalUp}{seedStr}

\begin{closurestatus}{\CofinalCompactIntervalUp}
  \constructivestory{A $\CofinalCompactIntervalUp$ packet is generated from a ClosedBoundedIntervalUp source row, a CompactIntervalLocatorUp row, a DyadicRatCoreUp tolerance ledger, StreamNameUp windows, RegSeqRatUp readback, a RealUp seal, componentwise hsame transport, Cont replay, Pkg provenance, and a local NameCert row.}
  \theoryclosure{\seedClosure}
  \scopeclosed{\autoref{def:cofinal-compact-interval-carrier}, \autoref{thm:cofinal-compact-interval-namecert-obligations}, and \autoref{thm:cofinal-compact-interval-real-handoff} seed the finite cofinal compact-interval route from closed bounded intervals through locator and dyadic-window rows to a Real-facing seal.}
  \formalstatus{\unformalizedV}
  \bridgestatus{none}
  \notclaimed{This chapter does not close open-cover compactness, selected subcovers, quotient real equality, countable choice, trichotomy, host interval equality, ambient $\RealUp$ completeness, Lean verification, or bridge checking.}
  \upgradepath{Obligation closure requires checked carrier admission, interval-source exactness, locator cofinality, finite-window routing, regular readback, Real-seal handoff, structural transport, replay, provenance, and local naming for BEDC.Derived.CofinalCompactIntervalUp.}
\end{closurestatus}
